% Môn: Vật lý
% Lớp: 10
% Chương: Chương II. Động học
% Bài: L10C2 Bài 4. Độ dịch chuyển và quãng đường đi được
% Loại: Lý thuyết
% File riêng, không trộn vào de.tex
% Định nghĩa lệnh Lý thuyết — dùng khi biên dịch lt.tex bằng TeX.
% App web đọc lt.tex trực tiếp (bỏ \input file này).
% Trong lt.tex, từ thư mục bài:
%   % Định nghĩa lệnh Lý thuyết — dùng khi biên dịch lt.tex bằng TeX.
% App web đọc lt.tex trực tiếp (bỏ \input file này).
% Trong lt.tex, từ thư mục bài:
%   % Định nghĩa lệnh Lý thuyết — dùng khi biên dịch lt.tex bằng TeX.
% App web đọc lt.tex trực tiếp (bỏ \input file này).
% Trong lt.tex, từ thư mục bài:
%   \input{../../../../_lenh/lythuyet.tex}
% từ gốc repo:
%   \input{ngan-hang/_lenh/lythuyet.tex}
\makeatletter
\providecommand{\indam}[1]{\textbf{#1}}
\providecommand{\chuy}[1]{\par\smallskip\noindent\textit{#1}\par}
\providecommand{\luuy}[1]{\par\smallskip\noindent\textbf{Lưu ý.} #1\par}
\providecommand{\ghichu}[1]{\par\smallskip\noindent\textbf{Ghi chú.} #1\par}
\providecommand{\vidu}[1]{\par\smallskip\noindent\textbf{Ví dụ.} #1\par}
\providecommand{\Blythuyet}{\subsection*{Lý thuyết}}
% Hai cột: kiến thức | hình
\providecommand{\haicotchay}[2]{%
  \par\medskip\noindent
  \begin{minipage}[t]{0.52\linewidth}#1\end{minipage}\hfill
  \begin{minipage}[t]{0.46\linewidth}#2\end{minipage}%
  \par\medskip
}
\providecommand{\kienthuccot}[2]{\haicotchay{#1}{#2}}
% Dự phòng nếu chưa nạp graphicx / caption (không ghi đè bản thật)
\providecommand{\resizebox}[3]{#3}
\providecommand{\captionof}[2]{\par\smallskip\noindent\textit{#2}\par}
\newenvironment{hoatdong}[1]{%
  \par\medskip\noindent\textbf{Hoạt động.} #1\par\smallskip
}{\par\medskip}
\newenvironment{emcobiet}{%
  \par\medskip\noindent\textbf{Em có biết}\par\smallskip
}{\par\medskip}
\newenvironment{emdahoc}{%
  \par\medskip\noindent\textbf{Em đã học}\par\smallskip
}{\par\medskip}
\newenvironment{emcothe}{%
  \par\medskip\noindent\textbf{Em có thể}\par\smallskip
}{\par\medskip}
\newenvironment{kienthuc}{%
  \par\medskip\noindent\textbf{Kiến thức cốt lõi}\par\smallskip
}{\par\medskip}
\newenvironment{traloi}{%
  \par\smallskip\noindent\textit{Trả lời / ghi chú của em}\par
  \vspace{1.2em}%
}{\par\medskip}
\makeatother

% từ gốc repo:
%   % Định nghĩa lệnh Lý thuyết — dùng khi biên dịch lt.tex bằng TeX.
% App web đọc lt.tex trực tiếp (bỏ \input file này).
% Trong lt.tex, từ thư mục bài:
%   \input{../../../../_lenh/lythuyet.tex}
% từ gốc repo:
%   \input{ngan-hang/_lenh/lythuyet.tex}
\makeatletter
\providecommand{\indam}[1]{\textbf{#1}}
\providecommand{\chuy}[1]{\par\smallskip\noindent\textit{#1}\par}
\providecommand{\luuy}[1]{\par\smallskip\noindent\textbf{Lưu ý.} #1\par}
\providecommand{\ghichu}[1]{\par\smallskip\noindent\textbf{Ghi chú.} #1\par}
\providecommand{\vidu}[1]{\par\smallskip\noindent\textbf{Ví dụ.} #1\par}
\providecommand{\Blythuyet}{\subsection*{Lý thuyết}}
% Hai cột: kiến thức | hình
\providecommand{\haicotchay}[2]{%
  \par\medskip\noindent
  \begin{minipage}[t]{0.52\linewidth}#1\end{minipage}\hfill
  \begin{minipage}[t]{0.46\linewidth}#2\end{minipage}%
  \par\medskip
}
\providecommand{\kienthuccot}[2]{\haicotchay{#1}{#2}}
% Dự phòng nếu chưa nạp graphicx / caption (không ghi đè bản thật)
\providecommand{\resizebox}[3]{#3}
\providecommand{\captionof}[2]{\par\smallskip\noindent\textit{#2}\par}
\newenvironment{hoatdong}[1]{%
  \par\medskip\noindent\textbf{Hoạt động.} #1\par\smallskip
}{\par\medskip}
\newenvironment{emcobiet}{%
  \par\medskip\noindent\textbf{Em có biết}\par\smallskip
}{\par\medskip}
\newenvironment{emdahoc}{%
  \par\medskip\noindent\textbf{Em đã học}\par\smallskip
}{\par\medskip}
\newenvironment{emcothe}{%
  \par\medskip\noindent\textbf{Em có thể}\par\smallskip
}{\par\medskip}
\newenvironment{kienthuc}{%
  \par\medskip\noindent\textbf{Kiến thức cốt lõi}\par\smallskip
}{\par\medskip}
\newenvironment{traloi}{%
  \par\smallskip\noindent\textit{Trả lời / ghi chú của em}\par
  \vspace{1.2em}%
}{\par\medskip}
\makeatother

\makeatletter
\providecommand{\indam}[1]{\textbf{#1}}
\providecommand{\chuy}[1]{\par\smallskip\noindent\textit{#1}\par}
\providecommand{\luuy}[1]{\par\smallskip\noindent\textbf{Lưu ý.} #1\par}
\providecommand{\ghichu}[1]{\par\smallskip\noindent\textbf{Ghi chú.} #1\par}
\providecommand{\vidu}[1]{\par\smallskip\noindent\textbf{Ví dụ.} #1\par}
\providecommand{\Blythuyet}{\subsection*{Lý thuyết}}
% Hai cột: kiến thức | hình
\providecommand{\haicotchay}[2]{%
  \par\medskip\noindent
  \begin{minipage}[t]{0.52\linewidth}#1\end{minipage}\hfill
  \begin{minipage}[t]{0.46\linewidth}#2\end{minipage}%
  \par\medskip
}
\providecommand{\kienthuccot}[2]{\haicotchay{#1}{#2}}
% Dự phòng nếu chưa nạp graphicx / caption (không ghi đè bản thật)
\providecommand{\resizebox}[3]{#3}
\providecommand{\captionof}[2]{\par\smallskip\noindent\textit{#2}\par}
\newenvironment{hoatdong}[1]{%
  \par\medskip\noindent\textbf{Hoạt động.} #1\par\smallskip
}{\par\medskip}
\newenvironment{emcobiet}{%
  \par\medskip\noindent\textbf{Em có biết}\par\smallskip
}{\par\medskip}
\newenvironment{emdahoc}{%
  \par\medskip\noindent\textbf{Em đã học}\par\smallskip
}{\par\medskip}
\newenvironment{emcothe}{%
  \par\medskip\noindent\textbf{Em có thể}\par\smallskip
}{\par\medskip}
\newenvironment{kienthuc}{%
  \par\medskip\noindent\textbf{Kiến thức cốt lõi}\par\smallskip
}{\par\medskip}
\newenvironment{traloi}{%
  \par\smallskip\noindent\textit{Trả lời / ghi chú của em}\par
  \vspace{1.2em}%
}{\par\medskip}
\makeatother

% từ gốc repo:
%   % Định nghĩa lệnh Lý thuyết — dùng khi biên dịch lt.tex bằng TeX.
% App web đọc lt.tex trực tiếp (bỏ \input file này).
% Trong lt.tex, từ thư mục bài:
%   % Định nghĩa lệnh Lý thuyết — dùng khi biên dịch lt.tex bằng TeX.
% App web đọc lt.tex trực tiếp (bỏ \input file này).
% Trong lt.tex, từ thư mục bài:
%   \input{../../../../_lenh/lythuyet.tex}
% từ gốc repo:
%   \input{ngan-hang/_lenh/lythuyet.tex}
\makeatletter
\providecommand{\indam}[1]{\textbf{#1}}
\providecommand{\chuy}[1]{\par\smallskip\noindent\textit{#1}\par}
\providecommand{\luuy}[1]{\par\smallskip\noindent\textbf{Lưu ý.} #1\par}
\providecommand{\ghichu}[1]{\par\smallskip\noindent\textbf{Ghi chú.} #1\par}
\providecommand{\vidu}[1]{\par\smallskip\noindent\textbf{Ví dụ.} #1\par}
\providecommand{\Blythuyet}{\subsection*{Lý thuyết}}
% Hai cột: kiến thức | hình
\providecommand{\haicotchay}[2]{%
  \par\medskip\noindent
  \begin{minipage}[t]{0.52\linewidth}#1\end{minipage}\hfill
  \begin{minipage}[t]{0.46\linewidth}#2\end{minipage}%
  \par\medskip
}
\providecommand{\kienthuccot}[2]{\haicotchay{#1}{#2}}
% Dự phòng nếu chưa nạp graphicx / caption (không ghi đè bản thật)
\providecommand{\resizebox}[3]{#3}
\providecommand{\captionof}[2]{\par\smallskip\noindent\textit{#2}\par}
\newenvironment{hoatdong}[1]{%
  \par\medskip\noindent\textbf{Hoạt động.} #1\par\smallskip
}{\par\medskip}
\newenvironment{emcobiet}{%
  \par\medskip\noindent\textbf{Em có biết}\par\smallskip
}{\par\medskip}
\newenvironment{emdahoc}{%
  \par\medskip\noindent\textbf{Em đã học}\par\smallskip
}{\par\medskip}
\newenvironment{emcothe}{%
  \par\medskip\noindent\textbf{Em có thể}\par\smallskip
}{\par\medskip}
\newenvironment{kienthuc}{%
  \par\medskip\noindent\textbf{Kiến thức cốt lõi}\par\smallskip
}{\par\medskip}
\newenvironment{traloi}{%
  \par\smallskip\noindent\textit{Trả lời / ghi chú của em}\par
  \vspace{1.2em}%
}{\par\medskip}
\makeatother

% từ gốc repo:
%   % Định nghĩa lệnh Lý thuyết — dùng khi biên dịch lt.tex bằng TeX.
% App web đọc lt.tex trực tiếp (bỏ \input file này).
% Trong lt.tex, từ thư mục bài:
%   \input{../../../../_lenh/lythuyet.tex}
% từ gốc repo:
%   \input{ngan-hang/_lenh/lythuyet.tex}
\makeatletter
\providecommand{\indam}[1]{\textbf{#1}}
\providecommand{\chuy}[1]{\par\smallskip\noindent\textit{#1}\par}
\providecommand{\luuy}[1]{\par\smallskip\noindent\textbf{Lưu ý.} #1\par}
\providecommand{\ghichu}[1]{\par\smallskip\noindent\textbf{Ghi chú.} #1\par}
\providecommand{\vidu}[1]{\par\smallskip\noindent\textbf{Ví dụ.} #1\par}
\providecommand{\Blythuyet}{\subsection*{Lý thuyết}}
% Hai cột: kiến thức | hình
\providecommand{\haicotchay}[2]{%
  \par\medskip\noindent
  \begin{minipage}[t]{0.52\linewidth}#1\end{minipage}\hfill
  \begin{minipage}[t]{0.46\linewidth}#2\end{minipage}%
  \par\medskip
}
\providecommand{\kienthuccot}[2]{\haicotchay{#1}{#2}}
% Dự phòng nếu chưa nạp graphicx / caption (không ghi đè bản thật)
\providecommand{\resizebox}[3]{#3}
\providecommand{\captionof}[2]{\par\smallskip\noindent\textit{#2}\par}
\newenvironment{hoatdong}[1]{%
  \par\medskip\noindent\textbf{Hoạt động.} #1\par\smallskip
}{\par\medskip}
\newenvironment{emcobiet}{%
  \par\medskip\noindent\textbf{Em có biết}\par\smallskip
}{\par\medskip}
\newenvironment{emdahoc}{%
  \par\medskip\noindent\textbf{Em đã học}\par\smallskip
}{\par\medskip}
\newenvironment{emcothe}{%
  \par\medskip\noindent\textbf{Em có thể}\par\smallskip
}{\par\medskip}
\newenvironment{kienthuc}{%
  \par\medskip\noindent\textbf{Kiến thức cốt lõi}\par\smallskip
}{\par\medskip}
\newenvironment{traloi}{%
  \par\smallskip\noindent\textit{Trả lời / ghi chú của em}\par
  \vspace{1.2em}%
}{\par\medskip}
\makeatother

\makeatletter
\providecommand{\indam}[1]{\textbf{#1}}
\providecommand{\chuy}[1]{\par\smallskip\noindent\textit{#1}\par}
\providecommand{\luuy}[1]{\par\smallskip\noindent\textbf{Lưu ý.} #1\par}
\providecommand{\ghichu}[1]{\par\smallskip\noindent\textbf{Ghi chú.} #1\par}
\providecommand{\vidu}[1]{\par\smallskip\noindent\textbf{Ví dụ.} #1\par}
\providecommand{\Blythuyet}{\subsection*{Lý thuyết}}
% Hai cột: kiến thức | hình
\providecommand{\haicotchay}[2]{%
  \par\medskip\noindent
  \begin{minipage}[t]{0.52\linewidth}#1\end{minipage}\hfill
  \begin{minipage}[t]{0.46\linewidth}#2\end{minipage}%
  \par\medskip
}
\providecommand{\kienthuccot}[2]{\haicotchay{#1}{#2}}
% Dự phòng nếu chưa nạp graphicx / caption (không ghi đè bản thật)
\providecommand{\resizebox}[3]{#3}
\providecommand{\captionof}[2]{\par\smallskip\noindent\textit{#2}\par}
\newenvironment{hoatdong}[1]{%
  \par\medskip\noindent\textbf{Hoạt động.} #1\par\smallskip
}{\par\medskip}
\newenvironment{emcobiet}{%
  \par\medskip\noindent\textbf{Em có biết}\par\smallskip
}{\par\medskip}
\newenvironment{emdahoc}{%
  \par\medskip\noindent\textbf{Em đã học}\par\smallskip
}{\par\medskip}
\newenvironment{emcothe}{%
  \par\medskip\noindent\textbf{Em có thể}\par\smallskip
}{\par\medskip}
\newenvironment{kienthuc}{%
  \par\medskip\noindent\textbf{Kiến thức cốt lõi}\par\smallskip
}{\par\medskip}
\newenvironment{traloi}{%
  \par\smallskip\noindent\textit{Trả lời / ghi chú của em}\par
  \vspace{1.2em}%
}{\par\medskip}
\makeatother

\makeatletter
\providecommand{\indam}[1]{\textbf{#1}}
\providecommand{\chuy}[1]{\par\smallskip\noindent\textit{#1}\par}
\providecommand{\luuy}[1]{\par\smallskip\noindent\textbf{Lưu ý.} #1\par}
\providecommand{\ghichu}[1]{\par\smallskip\noindent\textbf{Ghi chú.} #1\par}
\providecommand{\vidu}[1]{\par\smallskip\noindent\textbf{Ví dụ.} #1\par}
\providecommand{\Blythuyet}{\subsection*{Lý thuyết}}
% Hai cột: kiến thức | hình
\providecommand{\haicotchay}[2]{%
  \par\medskip\noindent
  \begin{minipage}[t]{0.52\linewidth}#1\end{minipage}\hfill
  \begin{minipage}[t]{0.46\linewidth}#2\end{minipage}%
  \par\medskip
}
\providecommand{\kienthuccot}[2]{\haicotchay{#1}{#2}}
% Dự phòng nếu chưa nạp graphicx / caption (không ghi đè bản thật)
\providecommand{\resizebox}[3]{#3}
\providecommand{\captionof}[2]{\par\smallskip\noindent\textit{#2}\par}
\newenvironment{hoatdong}[1]{%
  \par\medskip\noindent\textbf{Hoạt động.} #1\par\smallskip
}{\par\medskip}
\newenvironment{emcobiet}{%
  \par\medskip\noindent\textbf{Em có biết}\par\smallskip
}{\par\medskip}
\newenvironment{emdahoc}{%
  \par\medskip\noindent\textbf{Em đã học}\par\smallskip
}{\par\medskip}
\newenvironment{emcothe}{%
  \par\medskip\noindent\textbf{Em có thể}\par\smallskip
}{\par\medskip}
\newenvironment{kienthuc}{%
  \par\medskip\noindent\textbf{Kiến thức cốt lõi}\par\smallskip
}{\par\medskip}
\newenvironment{traloi}{%
  \par\smallskip\noindent\textit{Trả lời / ghi chú của em}\par
  \vspace{1.2em}%
}{\par\medskip}
\makeatother


\subsection{Lý thuyết}

\begin{emdahoc}
\begin{itemize}
    \item Khái niệm về chuyển động cơ học, vật mốc và hệ quy chiếu.
    \item Khái niệm về tốc độ và vận tốc.
\end{itemize}
\end{emdahoc}

\subsubsection{1. Độ dịch chuyển}

\begin{kienthuc}
\begin{itemize}
    \item \textbf{Độ dịch chuyển} là một đại lượng vectơ, đặc trưng cho sự thay đổi vị trí của vật.
    \item Độ dịch chuyển được xác định bằng vectơ nối vị trí đầu và vị trí cuối của vật.
    \item Kí hiệu: $\vec{d}$.
\end{itemize}
\end{kienthuc}

\begin{hoatdong}
Một người đi bộ từ nhà (điểm O) đến siêu thị (điểm A) cách nhà $600 \text{m}$ theo hướng Đông, rồi quay lại đi về phía Tây $200 \text{m}$ đến điểm $B$ để mua thêm đồ ở cửa hàng.
\begin{itemize}
    \item Vẽ hình minh họa chuyển động của người đó.
    \item Xác định vị trí đầu và vị trí cuối của người đó.
    \item Vẽ vectơ độ dịch chuyển của người đó từ O đến A, từ A đến B và từ O đến B.
\end{itemize}
\end{hoatdong}

\begin{kienthuc}
\begin{itemize}
    \item Trong chuyển động thẳng, nếu chọn trục tọa độ trùng với quỹ đạo chuyển động, độ dịch chuyển có thể được biểu diễn bằng một giá trị đại số.
    \item Nếu vật chuyển động từ vị trí có tọa độ $x_1$ đến vị trí có tọa độ $x_2$, độ dịch chuyển là $\Delta x = x_2 - x_1$.
    \item Giá trị $\Delta x$ có thể dương, âm hoặc bằng không, tùy thuộc vào chiều chuyển động và vị trí tương đối của $x_1, x_2$.
    \item Độ lớn của độ dịch chuyển là $|\Delta x|$.
\end{itemize}
\end{kienthuc}

\subsubsection{2. Quãng đường đi được}

\begin{kienthuc}
\begin{itemize}
    \item \textbf{Quãng đường đi được} là tổng chiều dài quỹ đạo mà vật đã đi qua trong một khoảng thời gian nhất định.
    \item Quãng đường là một đại lượng vô hướng và luôn có giá trị không âm.
    \item Kí hiệu: $s$.
\end{itemize}
\end{kienthuc}

\begin{hoatdong}
Trở lại ví dụ người đi bộ từ nhà (O) đến siêu thị (A) cách nhà $600 \text{m}$ theo hướng Đông, rồi quay lại đi về phía Tây $200 \text{m}$ đến điểm $B$.
\begin{itemize}
    \item Tính quãng đường người đó đi được từ O đến A.
    \item Tính quãng đường người đó đi được từ A đến B.
    \item Tính tổng quãng đường người đó đi được từ O đến B.
\end{itemize}
\end{hoatdong}

\subsubsection{3. Phân biệt quãng đường và độ dịch chuyển}

\begin{kienthuc}
\begin{itemize}
    \item \textbf{Độ dịch chuyển} là đại lượng vectơ, chỉ phụ thuộc vào vị trí đầu và vị trí cuối.
    \item \textbf{Quãng đường} là đại lượng vô hướng, phụ thuộc vào toàn bộ quỹ đạo chuyển động.
    \item Trong chuyển động thẳng và không đổi chiều, độ lớn của độ dịch chuyển bằng quãng đường đi được ($|\Delta x| = s$).
    \item Trong các trường hợp khác, quãng đường thường lớn hơn độ lớn của độ dịch chuyển ($s \ge |\Delta x|$).
\end{itemize}
\end{kienthuc}

\begin{hoatdong}
Một người đi bộ từ nhà (điểm O) đến siêu thị (điểm A) cách nhà $600 \text{m}$ theo hướng Đông, rồi quay lại đi về phía Tây $200 \text{m}$ đến điểm $B$ để mua thêm đồ ở cửa hàng.
\begin{itemize}
    \item Chọn gốc tọa độ tại O, chiều dương là chiều từ O đến A (hướng Đông).
    \item Xác định tọa độ của A và B.
    \item Tính độ dịch chuyển của người đó khi đi từ O đến B.
    \item Tính quãng đường người đó đi được khi đi từ O đến B.
    \item So sánh độ lớn của độ dịch chuyển và quãng đường trong trường hợp này.
\end{itemize}
\end{hoatdong}

\subsubsection{4. Đồ thị độ dịch chuyển -- thời gian}

\begin{kienthuc}
\begin{itemize}
    \item Đồ thị độ dịch chuyển -- thời gian ($x-t$) biểu diễn sự phụ thuộc của độ dịch chuyển (hoặc tọa độ) của vật vào thời gian.
    \item Trục tung biểu diễn độ dịch chuyển (hoặc tọa độ $x$), trục hoành biểu diễn thời gian ($t$).
    \item Từ đồ thị $x-t$, ta có thể xác định:
    \begin{itemize}
        \item Vị trí của vật tại các thời điểm khác nhau.
        \item Chiều chuyển động của vật (độ dốc của đồ thị).
        \item Vận tốc của vật (độ dốc của đồ thị).
        \item Độ dịch chuyển trong một khoảng thời gian.
    \end{itemize}
\end{itemize}
\end{kienthuc}

\begin{haicotchay}{Đồ thị độ dịch chuyển -- thời gian}{
\begin{tikzpicture}
\draw[->] (0,0) -- (5,0) node[right] {$t (\text{s})$};
\draw[->] (0,0) -- (0,4) node[above] {$x (\text{m})$};
\draw[thick, blue] (0,1) -- (2,3) node[above right] {A};
\draw[thick, blue] (2,3) -- (4,1) node[below right] {B};
\node at (0,1) [left] {$x_0$};
\node at (2,3) [above left] {$x_1$};
\node at (4,1) [below left] {$x_2$};
\node at (2,0) [below] {$t_1$};
\node at (4,0) [below] {$t_2$};
\end{tikzpicture}
}
\end{haicotchay}

\begin{kienthuc}
\begin{itemize}
    \item Nếu đồ thị là đường thẳng dốc lên, vật chuyển động theo chiều dương.
    \item Nếu đồ thị là đường thẳng dốc xuống, vật chuyển động theo chiều âm.
    \item Nếu đồ thị là đường thẳng nằm ngang, vật đứng yên.
    \item Độ dốc của đường thẳng trên đồ thị $x-t$ cho biết vận tốc của vật: $v = \frac{\Delta x}{\Delta t}$.
\end{itemize}
\end{kienthuc}

\subsubsection{5. Tổng hợp độ dịch chuyển}

\begin{kienthuc}
\begin{itemize}
    \item Khi một vật thực hiện nhiều độ dịch chuyển liên tiếp, độ dịch chuyển tổng hợp là tổng vectơ của các độ dịch chuyển thành phần.
    \item Nếu vật dịch chuyển từ A đến B ($\vec{d}_{AB}$) rồi từ B đến C ($\vec{d}_{BC}$), thì độ dịch chuyển tổng hợp từ A đến C là $\vec{d}_{AC} = \vec{d}_{AB} + \vec{d}_{BC}$.
    \item Phép tổng hợp độ dịch chuyển tuân theo quy tắc cộng vectơ (quy tắc hình bình hành hoặc quy tắc đa giác).
\end{itemize}
\end{kienthuc}

\begin{hoatdong}
Hai người đi xe đạp từ $A$ đến $C$. Người thứ nhất đi theo đường từ $A$ đến $B$ rồi từ $B$ đến $C$; người thứ hai đi thẳng từ $A$ đến $C$. Cả hai đều về đích cùng một lúc.
\begin{itemize}
    \item Vẽ hình minh họa hai con đường đi.
    \item So sánh độ dịch chuyển của hai người.
    \item So sánh quãng đường đi được của hai người.
\end{itemize}
\end{hoatdong}

\begin{emcothe}
\begin{itemize}
    \item Vận dụng khái niệm độ dịch chuyển và quãng đường để giải thích các hiện tượng trong đời sống.
    \item Phân tích và vẽ đồ thị độ dịch chuyển -- thời gian cho các chuyển động đơn giản.
    \item Tính toán độ dịch chuyển tổng hợp bằng phương pháp vectơ trong các trường hợp cụ thể.
\end{itemize}
\end{emcothe}
